\section{Training Details}
\label{sec:supp-training}

\subsection{Data Sampling and Curriculum}

The second stage is initialized from the EMA weights of the SpatialCog
warm-start checkpoint, with optimizer and scheduler states reset.
We detail the source-sampling strategy, curriculum transitions, structured
text conditioning, and learning-rate schedule used during its 51,100
optimization steps.

\paragraph{Hierarchical source sampling.}
We organize the training sources into five sampling domains.
Simulation combines Sim7 and MolmoSpaces; Molmo-real comprises DROID,
EgoDex, Stereo4D, and Xperience; Outdoor contains YT-VIS; and Real3 and
SpatialCog retain their collection-level definitions.
These groups determine curriculum sampling and differ from the four
data collections described in the main paper.

At optimization step $k$, the probability of selecting source $s$
in domain $d$ is
\begin{equation}
p_k(s)
=
p_k(d)
\frac{N_s^{\alpha_d}}
{\sum_{s' \in d} N_{s'}^{\alpha_d}},
\qquad
\alpha_d =
\begin{cases}
0.8, & d = \mathrm{Real3}, \\
0.5, & \text{otherwise}.
\end{cases}
\label{eq:supp-source-sampling}
\end{equation}
Here, $p_k(d)$ is the curriculum-dependent domain probability, and $N_s$
is the source count used to compute the sampler's within-domain
probabilities.
The sublinear exponent reduces the dominance of larger sources within
each domain. Sim7 is treated as a single source and retains the natural
proportions of its constituent datasets.

\paragraph{Curriculum intervals.}
Table~\ref{tab:supp-curriculum} specifies five sampling intervals that
implement the three phases described in the main paper: simulation-focused
training, mixed real-world adaptation, and SpatialCog refinement.
Real-world adaptation first emphasizes Molmo-real, briefly introduces
Outdoor data, and then prioritizes Real3. Each interval retains
cross-domain replay.

We linearly interpolate adjacent domain-probability vectors over windows
centered at steps 12,000, 24,000, 24,100, and 48,100, with widths of
1,000, 50, 50, and 1,000 steps, respectively.
Outside these transition windows, domain probabilities remain fixed at
the values in Table~\ref{tab:supp-curriculum}.
With a global batch size of 128, the nominal second-stage budget is
6,540,800 sample draws, including repeated visits to the same clips.

\begin{table}[t]
\centering
\scriptsize
\setlength{\tabcolsep}{3pt}
\renewcommand{\arraystretch}{1.1}
\begin{tabular*}{\linewidth}
{@{\extracolsep{\fill}}lrrrrrr@{}}
\toprule
Interval
& Steps
& Simulation
& Molmo-real
& Outdoor
& Real3
& SpatialCog \\
\midrule
Simulation-focused
& $[0,12{,}000)$
& 80\% & 5\% & 0\% & 10\% & 5\% \\

Molmo-real adaptation
& $[12{,}000,24{,}000)$
& 10\% & 70\% & 0\% & 15\% & 5\% \\

Outdoor adaptation
& $[24{,}000,24{,}100)$
& 5\% & 20\% & 50\% & 25\% & 0\% \\

Real3 adaptation
& $[24{,}100,48{,}100)$
& 10\% & 10\% & 0\% & 75\% & 5\% \\

SpatialCog refinement
& $[48{,}100,51{,}100)$
& 5\% & 10\% & 0\% & 35\% & 50\% \\
\bottomrule
\end{tabular*}
\caption{Target domain-sampling probabilities for the second-stage
curriculum. Adjacent distributions are linearly interpolated near interval
boundaries using the windows specified in the text. Within each domain,
sources are sampled according to Eq.~\eqref{eq:supp-source-sampling}.}
\label{tab:supp-curriculum}
\end{table}

\paragraph{Structured text conditioning.}
Each prompt consists of the original task instruction followed by a
context block with four categorical fields.
Data origin distinguishes simulated from real-world observations.
Actor type is specified as robot, human, mixed, or none.
Camera context distinguishes egocentric, fixed external, and moving
external views. Environment is specified as indoor, outdoor, or mixed.
These attributes provide explicit context for interpreting the requested
motion across heterogeneous sources.

We use sample-level context metadata when available and predefined
source-level defaults otherwise. The context fields follow a fixed order:
origin, actor, view, and scene. Training and inference use the same prompt
format, which excludes dataset identifiers and flow-normalization values.

\paragraph{Learning-rate schedule.}
Figure~\ref{fig:supp-curriculum} shows the second-stage learning-rate
schedule. The learning rate increases linearly to $10^{-4}$ over the first
500 steps and remains at this level during the main simulation-focused
interval. Near the transition to real-world adaptation, it decreases to
$8\times10^{-5}$ and remains at this plateau until step 42,100, including
the brief Outdoor interval. It then decays to $2\times10^{-5}$ at step
48,100 and reaches $5\times10^{-6}$ at the end of SpatialCog refinement.

\begin{figure}[t]
\centering
\includegraphics[width=\linewidth]
{supplementary/figures/training_curriculum.pdf}
\caption{Learning-rate schedule for second-stage training.
Background bands indicate the curriculum intervals. The narrow gold band
marks the 100-step Outdoor interval, during which the learning rate remains
at $8\times10^{-5}$.}
\label{fig:supp-curriculum}
\end{figure}